% R008 — Guide to Cultivating Complex Analysis (Bahasa Indonesia)
% Reader-facing translation unit: source/ca-v1.9/ca.tex lines 936–959.
% Source release: v1.9 (commit a4d25d9ba37a486a5a020219eb8bd4e6602fd14c).
% Translation provenance: OpenAI Codex gpt-5.6-sol, Ultra, acting on the user's request.
% Preserve source glossary hooks, notation, and equation structure.

Lingkungan dasar (yaitu, bola terbuka) di $\C$
disebut \emph{\myindex{cakram}}.
Diberikan $p \in \C$ dan $r > 0$, definisikan cakram berjari-jari $r$ di
sekitar $p$ sebagai
\glsadd{not:disc}%
\begin{equation*}
\Delta_r(p)
\overset{\text{def}}{=}
\bigl\{ z \in \C : \sabs{z-p} < r \bigr\} .
\end{equation*}
Cakram yang berpusat di titik asal dan berjari-jari $1$ disebut
\emph{\myindex{cakram satuan}}
\glsadd{not:D}%
\begin{equation*}
\D
\overset{\text{def}}{=}
\Delta_1(0)
=
\bigl\{ z \in \C : \sabs{z} < 1 \bigr\} .
\end{equation*}
Cakram satuan akan sering muncul dalam kuliah ini karena ternyata banyak
bagian analisis kompleks dapat dikerjakan hanya dengan mempelajari cakram
satuan.
